% ================================================================
%  conteudo/semana02.tex
% ================================================================

\section{Semana 2 --- 23 de março de 2026 a 27 de março de 2026}

\begin{table}[H]
\centering
\caption{Registro de atividades da Semana 2.}
\begin{tabular}{llll}
\toprule
\textbf{Data} & \textbf{Dia} & \textbf{Horas} & \textbf{Descrição} \\
\midrule
23/03/2026 & Segunda-feira & 2,0 h & Atividades no CIA \\
24/03/2026 & Terça-feira & 2,0 h & Atividades no CIA \\
25/03/2026 & Quarta-feira & 4,0 h & Atividades no CIA \\
26/03/2026 & Quinta-feira & 4,0 h & Atividades no CIA \\
27/03/2026 & Sexta-feira & 5,5 h & Atividades no CIA \\
\midrule
\multicolumn{2}{l}{\textbf{Total da semana}} & \textbf{17,5 h} & \\
\multicolumn{2}{l}{\textbf{Total acumulado}} & \textbf{35,0 h} & \\
\bottomrule
\end{tabular}
\end{table}

\subsubsection*{Atividades realizadas}

Na segunda semana, as rotinas de análise e preparo de materiais foram
repetidas para todos os equipamentos do CIA, consolidando os procedimentos
aprendidos na semana anterior. Adicionalmente, foi desenvolvida uma
contribuição computacional relevante:

\begin{itemize}
    \item Continuidade das rotinas de análise e preparo de materiais para
          IRMS, TGA, RMN, BET e DSC;
    \item Manutenção e reabastecimento dos gases dos equipamentos;
    \item \textbf{Desenvolvimento de um \textit{script} em Python} para
          monitoramento e análise do nível de Hélio (He) e Nitrogênio
          ($\text{N}_2$) nos tanques criogênicos do laboratório;
    \item \textbf{Instalação do \textit{script}} no computador dedicado ao
          espectrômetro de RMN, integrando o monitoramento ao fluxo de
          trabalho do equipamento.
\end{itemize}

\subsubsection*{Observações}

O desenvolvimento do \textit{script} em Python representa uma contribuição
direta à infraestrutura do CIA, automatizando o acompanhamento dos níveis
de gases criogênicos e reduzindo o risco de interrupção das análises por
falta de suprimento. O RMN, por depender criticamente de hélio líquido para
a manutenção do campo magnético supercondutor, é o equipamento mais
sensível a variações no nível do criogênico.
